\section{Phương pháp}

\subsection{Giao thức xác minh theo cặp}

Thực nghiệm hiện tại dùng giao thức xác minh theo cặp (\textit{pair-based verification}). Mỗi mẫu đánh giá gồm một ảnh gallery và một ảnh probe. Nếu hai ảnh thuộc cùng danh tính, cặp được gọi là genuine pair; nếu khác danh tính, cặp được gọi là impostor pair. Hệ thống tính embedding cho hai ảnh và lấy cosine similarity:

\begin{equation}
    s(x_a, x_b) =
    \cos\left(\phi(x_a), \phi(x_b)\right)
    =
    \frac{\phi(x_a)^\top \phi(x_b)}
    {\|\phi(x_a)\|_2 \|\phi(x_b)\|_2}.
    \label{eq:cosine}
\end{equation}

Một genuine pair bị từ chối là false rejection; một impostor pair được chấp nhận là false acceptance. Giao thức này phù hợp để phân tích hành vi ngưỡng, FAR và FRR. Tuy nhiên, nó chưa phải là bài toán nhận diện phiên đầy đủ, nơi một probe được so với nhiều danh tính trong gallery và cần xử lý các tình huống open-set.

\subsection{Trích xuất ngữ cảnh môi trường}

Với mỗi ảnh đầu vào $x$, hệ thống tính vector ngữ cảnh:

\begin{equation}
    C(x) = [L(x), N(x), q(x)],
    \label{eq:context}
\end{equation}

\noindent trong đó $L$ là độ sáng chuẩn hóa, $N$ là nhiễu ước tính, và $q$ là độ tin cậy phát hiện hoặc proxy chất lượng. Độ sáng được tính từ kênh luminance:

\begin{equation}
    L = \frac{1}{HW} \sum_{i=1}^{H}\sum_{j=1}^{W}\frac{Y_{ij}}{255}.
    \label{eq:brightness}
\end{equation}

Nhiễu được ước tính từ phần dư sau làm mịn Gaussian:

\begin{equation}
    N = \frac{\operatorname{std}(Y - \operatorname{GaussianBlur}(Y))}{255}.
    \label{eq:noise}
\end{equation}

Giá trị $L$ được rời rạc hóa thành ba bin:

\begin{equation}
    b(L)=
    \begin{cases}
        \texttt{dark},   & L < 0{,}30,\\
        \texttt{medium}, & 0{,}30 \leq L \leq 0{,}60,\\
        \texttt{bright}, & L > 0{,}60.
    \end{cases}
    \label{eq:bins}
\end{equation}

Các bin này không được dùng để tự động hạ ngưỡng, mà dùng để hiệu chỉnh ngưỡng riêng theo phân bố calibration của từng điều kiện.

\subsection{Đường xử lý có điều kiện}

Hệ thống xét ba đường xử lý:

\begin{itemize}
    \item \textbf{Fast path}: không dùng enhancement bổ sung, ưu tiên độ trễ thấp cho ảnh đủ tốt.
    \item \textbf{Robust path}: dùng tiền xử lý nhẹ như CLAHE/gamma để tăng độ bền vững trong điều kiện xấu.
    \item \textbf{Defer path}: không tự động quyết định khi chất lượng hoặc score nằm trong vùng không chắc chắn.
\end{itemize}

Policy chọn đường xử lý được ký hiệu:

\begin{equation}
    r = \pi(C), \qquad r \in \{\texttt{fast}, \texttt{robust}, \texttt{defer}\}.
    \label{eq:path_policy}
\end{equation}

Robust preprocessing được xem là một đường xử lý ứng viên, không được giả định luôn cải thiện embedding. Vì vậy, báo cáo đánh giá robust path thông qua FRR/FAR/latency thay vì mặc định coi enhancement là tốt hơn.

\subsection{Ra quyết định có kiểm soát rủi ro}

Với FAR budget $\alpha$, ngưỡng risk-constrained được hiệu chỉnh trên tập calibration $\mathcal{D}_{cal}$ bằng bài toán:

\begin{equation}
    \tau_{\alpha,r,b}
    =
    \arg\min_{\tau}
    \operatorname{FRR}_{cal}(\tau \mid r,b)
    \quad
    \text{s.t.}
    \quad
    \operatorname{FAR}_{cal}(\tau \mid r,b) \leq \alpha.
    \label{eq:risk_threshold}
\end{equation}

Trong đó $r$ là đường xử lý, $b$ là bin điều kiện, và $\alpha \in \{0{,}01, 0{,}02, 0{,}03, 0{,}05\}$. Công thức~\ref{eq:risk_threshold} là điểm khác biệt chính so với ngưỡng thích nghi đơn thuần: ngưỡng chỉ được chọn nếu thỏa mãn ràng buộc FAR trên tập calibration.

Với phương pháp \method{M5}, hệ thống quyết định:

\begin{equation}
    \hat{y} =
    \begin{cases}
        \texttt{accept}, & s \geq \tau_{\alpha,r,b},\\
        \texttt{reject}, & s < \tau_{\alpha,r,b}.
    \end{cases}
    \label{eq:m5_decision}
\end{equation}

Với phương pháp \method{M6}, hệ thống thêm vùng defer quanh ngưỡng chấp nhận:

\begin{equation}
    \tau_{\text{accept}} = \tau_{\alpha,r,b},
    \qquad
    \tau_{\text{reject}} = \tau_{\text{accept}} - \delta,
    \label{eq:defer_thresholds}
\end{equation}

\noindent với $\delta=0{,}03$ trong thực nghiệm. Quyết định cuối cùng là:

\begin{equation}
    \hat{y} =
    \begin{cases}
        \texttt{accept}, & s \geq \tau_{\text{accept}},\\
        \texttt{reject}, & s < \tau_{\text{reject}},\\
        \texttt{defer}, & \tau_{\text{reject}} \leq s < \tau_{\text{accept}}.
    \end{cases}
    \label{eq:m6_decision}
\end{equation}

Nhờ defer, một mẫu không đủ chắc chắn không bị tự động chấp nhận chỉ vì môi trường xấu. Đây là cơ chế quan trọng để trả lời rủi ro an ninh của bài toán.

\subsection{Các phương pháp so sánh}

Bảng~\ref{tab:methods} tóm tắt các phương pháp được đánh giá trong benchmark AWS cuối cùng.

\begin{table}[!htbp]
\centering
\caption{Các phương pháp được so sánh trong giao thức risk-constrained}
\label{tab:methods}
\begin{tabularx}{\textwidth}{llX}
\toprule
\textbf{ID} & \textbf{Tên} & \textbf{Mô tả} \\
\midrule
\method{M0} & Always-fast fixed & Fast path với một ngưỡng cố định cho mọi điều kiện. \\
\method{M1} & Always-fast bin & Fast path với ngưỡng riêng cho từng bin sáng/trung bình/tối. \\
\method{M4} & Conditional path bin & Chọn fast/robust theo điều kiện, sau đó dùng ngưỡng theo path và bin. \\
\method{M5} & Conditional risk & Chọn path theo điều kiện và hiệu chỉnh ngưỡng sao cho FAR calibration $\leq \alpha$. \\
\method{M6} & Conditional risk + defer & Mở rộng \method{M5} bằng vùng defer giữa reject và accept. \\
\bottomrule
\end{tabularx}
\end{table}

\subsection{Metric đánh giá}

Ngoài FAR và FRR thông thường, báo cáo dùng các metric sau:

\begin{itemize}
    \item \textbf{FAR\_active}: FAR trên các mẫu mà hệ thống tự động accept/reject, không tính mẫu defer.
    \item \textbf{FRR\_active}: FRR trên các mẫu mà hệ thống tự động accept/reject.
    \item \textbf{FRR\_with\_defer\_as\_failure\_for\_genuine}: tính genuine bị defer như một trường hợp chưa phục vụ thành công.
    \item \textbf{defer\_rate}: tỷ lệ mẫu cần xác minh thêm.
    \item \textbf{automation\_rate}: $1-\text{defer\_rate}$, tỷ lệ mẫu được hệ thống xử lý tự động.
    \item \textbf{latency\_mean, latency\_p95, ram\_peak\_mb}: chi phí tài nguyên trong mô phỏng thiết bị biên.
\end{itemize}

Khi so sánh các phương pháp risk-constrained, kết quả chỉ có ý nghĩa nếu FAR\_active được kiểm tra cùng với FRR\_active và defer\_rate. Một phương pháp defer nhiều có thể giảm lỗi tự động, nhưng automation rate thấp sẽ làm giảm giá trị vận hành.
